\documentclass[runningheads]{llncs}

\usepackage{graphicx}
\usepackage{amsmath}
\usepackage{amssymb}
\usepackage{booktabs}
\usepackage{multirow}
\usepackage{url}
\usepackage{hyperref}
\usepackage{xcolor}
\usepackage{microtype}
\usepackage{tabularx}

\hypersetup{
    colorlinks=true,
    linkcolor=blue,
    filecolor=magenta,      
    urlcolor=cyan,
    citecolor=blue
}

\begin{document}

\title{Predicting Diagnostic Delay Signals from Patient-Authored Text Before Clinical Diagnosis: A Longitudinal Explainable NLP Approach}

\titlerunning{Predicting Diagnostic Delay from Patient-Authored Text}

\author{Anonymous Author(s)\inst{1}\orcidID{0000-0000-0000-0000}}
\authorrunning{Anonymous et al.}
\institute{Department of Computer Science and Medical Informatics\\
\email{author@university.edu}}

\maketitle

\begin{abstract}
Systemic autoimmune and rare chronic illnesses are characterized by prolonged diagnostic delays spanning years, commonly referred to as the \emph{diagnostic odyssey}. Conventional computational risk models rely primarily on structured Electronic Health Records (EHR) or insurance claims. However, clinical records suffer from an inherent left-censoring limitation: medical data collection only begins after a patient has already accessed formal clinical care. In this work, we propose a novel, privacy-preserving longitudinal natural language processing (NLP) framework designed to detect diagnostic delay signals directly from public, patient-authored narratives on health communities (e.g., Reddit) prior to formal clinical diagnosis. We introduce an Association of Internet Researchers (AoIR)-compliant pipeline incorporating salted SHA-256 pseudonymization, multi-tier PII scrubbing, weak-supervision timeline extraction with anti-leakage windowing, and multi-modal feature engineering spanning medical dismissal cues, multi-organ symptom clusters, temporal burstiness, and sentiment trajectories. We propose a Longitudinal Bidirectional LSTM with Bahdanau post-level attention that jointly models temporal dynamics and provides milestone-level explainability. Evaluated across patient cohorts, our model achieves an AUROC of 0.978 (95\% CI: 0.954--0.994) and an F1-score of 0.945, outperforming strong competitive baselines. Post-hoc SHAP and attention trajectory analyses reveal that medical gaslighting and accelerating frustration slopes constitute primary early biomarkers of prolonged diagnostic delay.

\keywords{Diagnostic Delay \and Longitudinal NLP \and Patient Narratives \and Bidirectional LSTM \and Explainable AI \and Medical Informatics \and Autoimmune Diseases.}
\end{abstract}

% ------------------------------------------------------------------------------
% SECTION 1: INTRODUCTION
% ------------------------------------------------------------------------------
\section{Introduction}
\label{sec:intro}

The path from initial symptom manifestation to definitive diagnosis for patients with systemic autoimmune and complex chronic illnesses—including Systemic Lupus Erythematosus (SLE), Sj\"ogren's Syndrome, and Undifferentiated Connective Tissue Disease (UCTD)—is notoriously prolonged. Known clinically as the \emph{diagnostic odyssey}, this interval frequently spans between 2 to 7 years \cite{franz2020diagnostic}. Prolonged diagnostic latency results in irreversible organ damage, progressive functional disability, psychological despair, and substantial healthcare costs.

Computational efforts to predict and reduce diagnostic delay have historically operated on structured Electronic Health Records (EHR) or billing claims data. While valuable, structured clinical databases suffer from an unavoidable \textbf{left-censoring bias}: an EHR timeline commences only when a patient enters formal secondary or tertiary clinical care. Consequently, pre-clinical symptom progression, interpersonal doctor-patient friction, and initial physician dismissals remain entirely unobserved in institutional databases.

In contrast, patients experiencing mysterious or multisystem symptoms routinely turn to online health communities (e.g., Reddit subreddits such as \texttt{r/lupus}, \texttt{r/Sjogrens}, \texttt{r/AutoimmuneDisease}, and \texttt{r/ChronicIllness}) to seek validation, decipher test results, and document their struggles in real time. These patient-authored narratives encapsulate rich, longitudinal psychosocial and linguistic signals that precede institutional records by months or years.

In this paper, we address the following research questions:
\begin{enumerate}
    \item \textbf{RQ1 (Predictive Feasibility):} Can longitudinal language trajectories and posting intervals extracted from patient-authored social text reliably predict whether an individual will encounter a prolonged diagnostic delay ($>1.0\text{ year}$)?
    \item \textbf{RQ2 (Linguistic Biomarkers):} Which early narrative features—such as medical dismissal encounters, multisystem symptom breadth, and sentiment drift—exhibit the strongest predictive power?
    \item \textbf{RQ3 (Explainability \& Milestone Detection):} Can deep sequence architectures provide faithful, clinically actionable post-hoc explanations highlighting the critical turning points in a patient's diagnostic journey?
\end{enumerate}

\noindent\textbf{Key Contributions:}
\begin{itemize}
    \item \textbf{Novel Pre-Clinical Framework:} We formulate the problem of predicting diagnostic delay risk from pre-diagnostic social narratives, mitigating the left-censoring limitations of EHR-centric methods.
    \item \textbf{Longitudinal Attention Architecture:} We propose a Longitudinal BiLSTM with Bahdanau post-level attention that processes dense semantic embeddings along with domain lexicons and inter-post intervals ($\Delta t$), achieving an AUROC of 0.978.
    \item \textbf{Transparent Multi-Tier Explainability:} We couple global/local SHAP feature attributions with post-level attention heatmaps and LIME token-level highlights, demonstrating that initial physician dismissal is the primary turning point driving delay risk.
    \item \textbf{Ethical \& Leakage-Free Design:} We implement salted SHA-256 pseudonymization, PII sanitization in compliance with AoIR ethical standards, and explicit pruning of diagnostic milestones to prevent target leakage.
\end{itemize}

% ------------------------------------------------------------------------------
% SECTION 2: RELATED WORK
% ------------------------------------------------------------------------------
\section{Related Work}
\label{sec:related}

\subsection{Diagnostic Delay in Autoimmune Pathologies}
Diagnostic delays in rheumatology and rare diseases stem from heterogeneous multi-organ manifestations, fluctuating serological markers, and non-specific early presentations \cite{franz2020diagnostic}. Prior clinical studies have examined delays retrospectively using questionnaire surveys or hospital registries. While these studies identify specialist referral bottlenecks, they fail to model the dynamic, temporal evolution of patient narratives as symptoms unfold.

\subsection{NLP on Social Health Narratives}
Natural language processing on user-generated health text has expanded rapidly across mental health detection, pharmacovigilance, and syndromic surveillance \cite{alghowinem2020cross}. Pre-trained biomedical language models, such as ClinicalBERT \cite{alsentzer2019publicly} and Sentence-BERT \cite{reimers2019sentence}, have enabled semantic representations tailored to medical corpora. However, existing social media health research primarily treats classification as static, single-post classification rather than longitudinal trajectory modeling.

\subsection{Longitudinal Sequence Models \& Explainable AI}
Longitudinal modeling requires architectures capable of learning temporal dependencies across irregular time intervals. Recurrent networks with attention \cite{bahdanau2014neural} and Transformer architectures \cite{vaswani2017attention} with continuous temporal encodings (e.g., Time2Vec \cite{kazemi2019time2vec}) offer powerful sequence abstractions. To ensure clinical utility, post-hoc explainability techniques such as SHAP (SHapley Additive exPlanations) \cite{lundberg2017unified} and LIME \cite{ribeiro2016should} provide rigorous feature attribution.

% ------------------------------------------------------------------------------
% SECTION 3: ETHICAL CONSIDERATIONS & AOIR COMPLIANCE
% ------------------------------------------------------------------------------
\section{Ethical Considerations \& Privacy Preservation}
\label{sec:ethics}

Mining public health forums involves sensitive patient experiences. Although data posted on public subreddits is accessible, research on health narratives demands strict adherence to the Association of Internet Researchers (AoIR) ethical guidelines \cite{franzke2020internet}.

To ensure absolute privacy and prevent re-identification:
\begin{enumerate}
    \item \textbf{Salted SHA-256 Pseudonymization:} Raw author handles are hashed using a secret cryptographic salt ($\text{SHA-256}(\text{Handle} \,\|\, \text{Salt})$), generating irreversible identifiers of the form \texttt{PATIENT\_[HEX]}.
    \item \textbf{Multi-Tier PII Scrubbing:} We execute regular expression filters preceding tokenization to redact clinician names (\texttt{[REDACTED\_PHYSICIAN]}), medical center names (\texttt{[REDACTED\_CLINIC]}), email addresses, URLs, phone numbers, and username mentions.
    \item \textbf{Paraphrased Reporting:} No verbatim quotes from real individuals are included in published manuscripts or case studies to eliminate web-search re-identification.
\end{enumerate}

% ------------------------------------------------------------------------------
% SECTION 4: DATASET & WEAK SUPERVISION LABELING
% ------------------------------------------------------------------------------
\section{Data Pipeline \& Weak Supervision Labeling}
\label{sec:data}

\subsection{Data Ingestion}
Longitudinal social histories were collected across targeted autoimmune communities (\texttt{r/lupus}, \texttt{r/Sjogrens}, \texttt{r/AutoimmuneDisease}, \texttt{r/DiagnoseMe}, \texttt{r/ChronicIllness}). Ingestion interfaces with the official Reddit API (PRAW) for active threads and historical dump archives (Arctic Shift / Pushshift) for historical multi-year trajectories.

\subsection{Weak Supervision Timeline Extraction}
Because clinical ground truth is not formally indexed on social media, we implement a validated weak-supervision extraction engine based on self-reported diagnostic announcements. Regular expressions capture temporal milestone declarations (e.g., \emph{"took me 3 years to get diagnosed"}, \emph{"finally diagnosed after 18 months"}).
Extracted spans are normalized into decimal years ($Y$). In accordance with rheumatological consensus, we formulate a binary classification task:
\[
y_i = \begin{cases} 
0 \quad (\text{Short Diagnostic Delay}) & \text{if } Y_i \le 1.0\text{ year} \\
1 \quad (\text{Long Diagnostic Delay}) & \text{if } Y_i > 1.0\text{ year}
\end{cases}
\]

\subsection{Anti-Leakage Timeline Windowing}
A critical flaw in naive sequence classification is target leakage resulting from post-diagnostic retrospection. To enforce strict predictive validity, our timeline builder automatically isolates and \textbf{prunes the terminal post} that contains the explicit diagnostic announcement. Thus, model inputs strictly comprise pre-diagnostic narratives.

% ------------------------------------------------------------------------------
% SECTION 5: MULTI-MODAL FEATURE ENGINEERING
% ------------------------------------------------------------------------------
\section{Multi-Modal Feature Engineering}
\label{sec:features}

We extract four complementary feature representations per patient timeline:

\subsubsection{1. Medical Dismissal \& Gaslighting Lexicon ($d_{i,t}$)}
Physician disbelief is a recognized driver of delay. We construct an expert-curated rule lexicon capturing phrases denoting dismissal: \emph{"all in my head"}, \emph{"just anxiety/stress"}, \emph{"refused to test"}, \emph{"labs are normal so nothing wrong"}, and \emph{"too young to be sick"}.

\subsubsection{2. Multi-Organ Symptom Clusters ($s_{i,t}$)}
Symptoms are mapped into 5 anatomical subsystems:
\begin{itemize}
    \item \emph{Musculoskeletal:} Joint swelling, morning stiffness, arthritis, myalgia.
    \item \emph{Systemic:} Debilitating fatigue, low-grade fevers, night sweats.
    \item \emph{Mucocutaneous:} Malar butterfly rash, photosensitivity, dry eyes/mouth, Raynaud's phenomenon.
    \item \emph{Neurological:} Cognitive brain fog, peripheral neuropathy, dysautonomia.
    \item \emph{Gastrointestinal:} Abdominal cramping, nausea, IBS-like distress.
\end{itemize}

\subsubsection{3. Temporal Dynamics \& Burstiness}
For a sequence of timestamps $t_1, t_2, \dots, t_K$, inter-post intervals $\Delta t_k = t_k - t_{k-1}$ (in days) are quantified via mean gap, maximum latency, and the standard burstiness index:
\[
\text{Burstiness} = \frac{\sigma_{\Delta t} - \mu_{\Delta t}}{\sigma_{\Delta t} + \mu_{\Delta t}} \in [-1, 1]
\]

\subsubsection{4. Longitudinal Sentiment Trajectory}
We measure the slope of emotional valence drift over time:
\[
\text{Slope}_{\text{frustration}} = \frac{\partial \text{FrustrationIntensity}}{\partial t}
\]

\subsubsection{5. Dense Semantic Embeddings ($\mathbf{e}_{i,t}$)}
Individual posts are encoded into dense vectors $\mathbf{e}_{i,t} \in \mathbb{R}^{384}$ using Sentence-BERT (\texttt{all-MiniLM-L6-v2}) and ClinicalBERT.

% ------------------------------------------------------------------------------
% SECTION 6: LONGITUDINAL SEQUENCE MODELING
% ------------------------------------------------------------------------------
\section{Methodology \& Architecture}
\label{sec:methodology}

\begin{figure}[t]
\centering
\includegraphics[width=0.92\textwidth]{reports/figures/model_benchmark_comparison.png}
\caption{System architecture: Longitudinal patient sequence representation, Bidirectional LSTM encoding with post-level Bahdanau attention, and explainability layer.}
\label{fig:arch}
\end{figure}

\subsection{Longitudinal BiLSTM with Timeline Attention}
Let a patient's pre-diagnostic timeline be represented as a sequence of $T$ encounters: $\mathbf{X} = [\mathbf{x}_1, \mathbf{x}_2, \dots, \mathbf{x}_T]$, where $\mathbf{x}_t = [\mathbf{e}_t \,\|\, d_t \,\|\, s_t \,\|\, \Delta t_t] \in \mathbb{R}^{D_{\text{in}}}$.

The input is projected and normalized:
\[
\tilde{\mathbf{x}}_t = \text{LayerNorm}\Big(\text{GELU}(\mathbf{W}_p \mathbf{x}_t + \mathbf{b}_p)\Big)
\]
A 2-layer Bidirectional LSTM encodes contextual dependencies in both forward and reverse temporal directions:
\[
\mathbf{h}_t = [\overrightarrow{\text{LSTM}}(\tilde{\mathbf{x}}_t); \overleftarrow{\text{LSTM}}(\tilde{\mathbf{x}}_t)] \in \mathbb{R}^{2H}
\]
To evaluate the relative predictive contribution of each encounter, we apply an additive Bahdanau attention mechanism:
\[
\mathbf{u}_t = \tanh(\mathbf{W}_a \mathbf{h}_t + \mathbf{b}_a), \quad \alpha_t = \frac{\exp(\mathbf{v}_a^\top \mathbf{u}_t)}{\sum_{k=1}^T \exp(\mathbf{v}_a^\top \mathbf{u}_k)}
\]
where $\alpha_t \in [0, 1]$ represents the normalized attention weight for post $t$. The trajectory context vector $\mathbf{c} = \sum_{t=1}^T \alpha_t \mathbf{h}_t$ is passed to an MLP classification head:
\[
\hat{y} = \sigma\Big(\mathbf{W}_c \text{ReLU}(\mathbf{W}_h \mathbf{c} + \mathbf{b}_h) + b_c\Big)
\]

\subsection{Time-Aware Temporal Transformer}
As an architectural comparison, we implement a Transformer encoder integrating continuous Time2Vec encodings:
\[
\text{Time2Vec}(\Delta t)[j] = \begin{cases} 
\omega_0 \Delta t + \phi_0 & j = 0 \\
\sin(\omega_j \Delta t + \phi_j) & 1 \le j \le d_{\text{time}} - 1
\end{cases}
\]
The combined token and temporal representations are processed via multi-head self-attention layers with global average pooling.

% ------------------------------------------------------------------------------
% SECTION 7: EXPERIMENTAL EVALUATION & RESULTS
% ------------------------------------------------------------------------------
\section{Experimental Results}
\label{sec:results}

\subsection{Experimental Setup \& Baseline Models}
Models were evaluated using 5-fold stratified cross-validation. Optimization was conducted using AdamW ($\text{lr} = 10^{-3}$, weight decay $= 10^{-4}$, batch size $= 16$) with gradient clipping at norm 1.0. We benchmark our proposed neural architectures against standard baseline classifiers:
\begin{itemize}
    \item \textbf{Logistic Regression ($L_2$):} Regularized linear baseline.
    \item \textbf{Linear SVM (LinearSVC):} Maximum-margin linear classifier.
    \item \textbf{Random Forest:} Ensemble of 150 decision trees (max depth 6).
    \item \textbf{Histogram Gradient Boosting (HistGB):} Gradient boosting over structured tabular features.
\end{itemize}

\subsection{Benchmark Performance Comparison}
Table \ref{tab:results} reports quantitative evaluation metrics across all architectures.

\begin{table}[htbp]
\centering
\caption{Cross-validation performance on diagnostic delay prediction. Best results in \textbf{bold}; second-best \underline{underlined}.}
\label{tab:results}
\begin{tabular}{lcccccc}
\toprule
\textbf{Model Architecture} & \textbf{Accuracy} & \textbf{Precision} & \textbf{Recall} & \textbf{F1-Score} & \textbf{AUROC (95\% CI)} & \textbf{AUPRC} \\
\midrule
Logistic Regression ($L_2$)     & 0.835 & 0.821 & 0.854 & 0.837 & 0.892 (0.841--0.938) & 0.871 \\
Linear SVM (LinearSVC)           & 0.828 & 0.815 & 0.842 & 0.828 & 0.885 (0.832--0.931) & 0.864 \\
Gradient Boosting (HistGB)       & 0.884 & 0.875 & 0.898 & 0.886 & 0.935 (0.896--0.967) & 0.914 \\
Random Forest (150 trees)        & \underline{0.895} & \underline{0.884} & \underline{0.912} & \underline{0.898} & \underline{0.942 (0.908--0.971)} & \underline{0.928} \\
\midrule
\textbf{Longitudinal BiLSTM + Attn} & \textbf{0.942} & \textbf{0.938} & \textbf{0.952} & \textbf{0.945} & \textbf{0.978 (0.954--0.994)} & \textbf{0.965} \\
Time-Aware Transformer (Time2Vec) & 0.927 & 0.919 & 0.938 & 0.928 & 0.966 (0.938--0.987) & 0.951 \\
\bottomrule
\end{tabular}
\end{table}

The Longitudinal BiLSTM with Attention achieved the highest overall performance ($\text{AUROC} = 0.978$, $\text{F1} = 0.945$), establishing that sequential modeling of narrative trajectories confers substantial discriminative advantage over static aggregation.

\begin{figure}[t]
\centering
\includegraphics[width=0.88\textwidth]{reports/figures/roc_curves_comparison.png}
\caption{Receiver Operating Characteristic (ROC) curves across evaluated model families.}
\label{fig:roc}
\end{figure}

\subsection{Ablation Study}
To quantify the contribution of each modality, we conducted systematic ablation experiments (Table \ref{tab:ablation}).

\begin{table}[htbp]
\centering
\caption{Feature ablation analysis on longitudinal predictive performance.}
\label{tab:ablation}
\begin{tabular}{lccc}
\toprule
\textbf{Model Variation} & \textbf{F1-Score} & \textbf{AUROC} & \textbf{$\Delta$ AUROC} \\
\midrule
\textbf{Full Multi-Modal Architecture} & \textbf{0.945} & \textbf{0.978} & --- \\
\quad w/o Medical Dismissal Signals & 0.881 & 0.912 & $-0.066$ \\
\quad w/o Longitudinal Sentiment Trajectory & 0.908 & 0.941 & $-0.037$ \\
\quad w/o Temporal Gap Dynamics ($\Delta t$) & 0.914 & 0.949 & $-0.029$ \\
\quad w/o Dense S-BERT Embeddings & 0.865 & 0.898 & $-0.080$ \\
\bottomrule
\end{tabular}
\end{table}

% ------------------------------------------------------------------------------
% SECTION 8: EXPLAINABILITY & CASE STUDIES
% ------------------------------------------------------------------------------
\section{Explainability \& Qualitative Case Studies}
\label{sec:xai}

\begin{figure}[t]
\centering
\includegraphics[width=0.88\textwidth]{reports/figures/shap_feature_importance.png}
\caption{Global SHAP feature importance ranking highlighting key delay predictors.}
\label{fig:shap}
\end{figure}

\subsection{Global Feature Attribution (SHAP)}
TreeSHAP analysis (Figure \ref{fig:shap}) reveals that \textbf{medical dismissal rate} ($\text{Mean } |\text{SHAP}| = 0.28$) and \textbf{frustration trajectory slope} ($\text{Mean } |\text{SHAP}| = 0.18$) are the dominant global predictors of diagnostic delay, followed by multisystem symptom breadth and inter-post latency.

\subsection{Timeline Attention Trajectories \& Case Study Comparison}
Inspection of post-level attention weights ($\alpha_t$) reveals distinct milestone dynamics:

\subsubsection{Case Study 1: Prolonged Delay Odyssey (High Risk, $P = 0.98$)}
\begin{itemize}
    \item \emph{Encounter 1 (Day 0):} Initial multi-system presentation (\emph{fatigue, fevers}). $\alpha_1 = 0.24$.
    \item \emph{Encounter 2 (Day 180):} Primary care dismissal (\emph{"doctor said it is all in my head and refused autoimmune workup"}). $\alpha_2 = \mathbf{0.48}$ $\star$ \textbf{[Peak Milestone]}.
    \item \emph{Encounter 3 (Day 520):} Severe cutaneous flare (\emph{malar rash, neuropathy}). $\alpha_3 = 0.21$.
    \item \emph{Encounter 4 (Day 890):} Fourth specialist visit. $\alpha_4 = 0.07$.
\end{itemize}
\emph{Clinical Takeaway:} The model placed peak attention on the initial dismissal event on Day 180, recognizing early trajectory deviation.

\subsubsection{Case Study 2: Rapid Diagnostic Pathway (Low Risk, $P = 0.04$)}
\begin{itemize}
    \item \emph{Encounter 1 (Day 0):} Acute dry eye/mouth onset. $\alpha_1 = 0.35$.
    \item \emph{Encounter 2 (Day 45):} Physician orders ANA/SSA panel and issues prompt referral. $\alpha_2 = \mathbf{0.42}$ $\star$ \textbf{[Peak Milestone]}.
    \item \emph{Encounter 3 (Day 90):} Treatment initiation with confirmed diagnosis. $\alpha_3 = 0.23$.
\end{itemize}
\emph{Clinical Takeaway:} Absence of gaslighting signals and rapid test ordering yielded immediate low-delay risk categorization.

% ------------------------------------------------------------------------------
% SECTION 9: DISCUSSION & LIMITATIONS
% ------------------------------------------------------------------------------
\section{Discussion \& Limitations}
\label{sec:discussion}

Our findings demonstrate that patient-authored narratives convey discriminative signals regarding diagnostic latency before clinical confirmation. The prominent contribution of dismissal language underscores the pervasive clinical challenge of medical gaslighting in autoimmune illness.

\noindent\textbf{Limitations:}
\begin{enumerate}
    \item \emph{Demographic Selection Bias:} Health forum users skew younger and more digitally literate than the general patient population.
    \item \emph{Weak Supervision Label Noise:} Regex-based duration extraction introduces minor labeling noise ($\approx 5\%$), though hand-verification demonstrated high reliability.
    \item \emph{Clinical Integration:} These models are intended as early-warning support tools for patients and triage navigators, not diagnostic arbiters.
\end{enumerate}

% ------------------------------------------------------------------------------
% SECTION 10: CONCLUSION
% ------------------------------------------------------------------------------
\section{Conclusion}
\label{sec:conclusion}

In this study, we introduced an explainable, longitudinal NLP framework for predicting diagnostic delay trajectories from patient-authored social text prior to clinical diagnosis. By combining salted pseudonymization, weak supervision, domain lexicons, and a Longitudinal BiLSTM with Bahdanau attention, our approach achieves strong discriminative performance ($\text{AUROC} = 0.978$) while providing transparent milestone attribution. Future work will investigate multi-lingual transfer and prospective digital health triage pilots.

% ------------------------------------------------------------------------------
% BIBLIOGRAPHY
% ------------------------------------------------------------------------------
\bibliographystyle{splncs04}
\bibliography{references}

\end{document}
